\section{Основные термины и определения}
\begin{itemize}
    \item \textbf{Стартап} --- это молодая, недавно созданная компания, находящаяся на стадии развития и строящая свой бизнес вокруг инновационной технологии, идеи или уникальной бизнес-модели.~\cite{vtbstartup}
    \item \textbf{Предиктор успеха} --- измеримый параметр или характеристика, статистически связанная с успехом или провалом стартапа.
    \item \textbf{Y Combinator} --- один из крупнейших и наиболее известных стартап-акселераторов в мире, основанный в 2005 году. Среди выпускников --- Airbnb, Dropbox, Stripe, Reddit и другие~\cite{ycombinator}.
    \item \textbf{Батч} --- набор компаний, получивших инвестиции в один период времени.
    \item \textbf{IPO} --- первичное публичное размещение акций на фондовой бирже, в результате которого компания становится публичной~\cite{ipo}.
    \item \textbf{Kaggle} --- публичная веб-платформа, на которой пользователи и организации могут публиковать наборы данных, исследовать и создавать модели, взаимодействовать с другими специалистами по данным.~\cite{wikikaggle}
    \item \textbf{Разведочный анализ данных (EDA)} --- анализ основных свойств данных, нахождение в них общих закономерностей, распределений и аномалий, зачастую с использованием инструментов визуализации.~\cite{wikiEDA}
    \item \textbf{Критерий Манна-Уитни} --- непараметрический статистический тест для сравнения двух независимых выборок, не требующий предположения о нормальности распределения.~\cite{mann-whitney}
    \item \textbf{Корреляция Пирсона} --- мера линейной зависимости между двумя переменными, принимающая значения от $-1$ до $+1$.~\cite{pearson-corr}
    \item \textbf{Ошибка выжившего} --- систематическая ошибка отбора, при которой анализируются только <<выжившие>> объекты, а <<не выжившие>> игнорируются, что приводит к искаженным выводам.
    \item \textbf{Когортный анализ} --- метод анализа, при котором объекты группируются по времени возникновения (когортам) и сравниваются внутри одной когорты, что позволяет контролировать временной фактор.
    \item \textbf{Crunchbase} --- это онлайн-платформа, которая предоставляет информацию о стартапах, компаниях, инвесторах и предпринимателях по всему миру.~\cite{getinvestor}
\end{itemize}
